% !TeX root = ../papers/nc-p3.tex
\section{四次超曲面としての全中心とPI次数}
\label{sec:quartic-center-order}
前節の開集合で折り畳み環を$B$、多項式部分環を$P=k[x_1,x_2,x_3,x_4]$、
残る中心四次式を$z$とする。以下で中心の変数を通常の次数1に数えるときは、
元の経路長を4で割っている。

\begin{theorem}\label{thm:quartic-center-order}
$t$を含む単一の非空開集合の全幾何点で、全中心は
\[
 Z(B)=P[z]/\bigl(z^4+c_1z^3+c_2z^2+c_3z+c_4\bigr),
       \qquad c_i\in P_i
\]
の表示を持つ（左の$z$は環の元、右は表示変数）。その射影三次元多様体は
幾何学的に整かつ正規な四次超曲面である。$B$はこの中心上有限な素環で、
中心の分数体上で次数16の中心的単純環となる。各頂点の隅環の次数は4である。
特にこの開集合のPI次数は16であり、成分の閉包の全標準AS点が$s_{32}=0$を満たす。
\end{theorem}
\begin{proof}
各隅$e_i B e_j$は$P$上自由階数64で、全体は自由階数1024である。
まず有理点$t$で、既出の四変数の整数係数リフトを取り、文字次数0に残る
第三の独立な四次式を$z$とする。順序と原始整数係数への規格化は
\path{computations/verify_quartic_center_order.py}に指定する。
$B/(x_1,\ldots,x_4)$の第0頂点の隅で$z$の乗法の最小多項式は
素数1009の剰余体で$T^4$である。従って剰余の$1,z,z^2,z^3$は独立である
（隅の単位は$e_0$とする）。全折り畳み環でも同じ四元は独立で、
剰余基底へ延長して斉次リフトを取ると$B$の$P$自由基底になる。
特にこの四元の張る部分加群は直和因子である。

中心変数を$(1,2,3,4)$に指定すると、同じ素数で四頂点の隅の直和に
矢との可換式を課した核は次元4となる。第0隅での$z$の最小多項式は
\[
 T^4+421T^3+601T^2+542T+727\pmod {1009}
\]
で重根を持たず、特性多項式はその十六乗である。
同じ隅の64基底の左正則作用を$L_u$とすると、
$(\operatorname{tr}(L_uL_v))_{u,v}$の階数は64である。

これらの有限計算の特殊化を説明する。定数$\alpha_i$を指定した商で、
固定始点・終点の長さが$n$以下かつ$n$と法4で合同な語を扱う。
一文字を先頭に付けた四組の前次数基底に、$4\mid n$のときは単位も加える。
六関係と$x_i-\alpha_i$に残りの長さの基底語を右から掛けて掃き出す。
これは中央変数$y$を経路長4として$x_i=\alpha_i y$とした斉次商の次数$n$と同じである。
元が$P$上自由なので、この商は$k[y]$上自由であり、その次数$n$の次元は
$(1+t+t^2+t^3)^4$の該当する剰余類の次数$n$までの係数和である。
実際の検算は、全始点で$n\le16$までこの値と一致することを確かめる。
従ってこれらの有限範囲の主小行列は混標数局所環でも単元に取れ、
剰余体の計算は標数零での階数下界を与える。第0隅の乗法は
長さ12以下の基底に長さ4の切断を左から掛けて長さ16の基底へ移して求める。
正則作用のtraceは、同じ終点を持つ四つの安定した64次元空間上の
矢行列の積から求める。

$K=\operatorname{Frac}P$とする。中心は四隅の対角元が十六の矢と可換である
という$P$上の有限行列の核である。上の一点の階数から$K$上で核次元は高々4、
既知の四元から丁度4となる。自由直和因子の飽和性により
\[
 Z(B)=B\cap\operatorname{span}_K(1,z,z^2,z^3)
      =P\oplus Pz\oplus Pz^2\oplus Pz^3 .
\]
従って$z^4$の一意な表示が定理のモニック関係を与え、斉次性から係数次数も従う。
この段階では中心の既約性を使っていない。

この四次の平面切断を検査する。$(x_1,x_2,x_3,x_4)=(s,u,s+2u,3s-u)$とし、
$s=1$, $u=0,1,2,3,4$で上の有限乗法を計算する。
どの剰余でも第0隅における$z$の最小多項式は四次である。
よって各値で$1,z,z^2,z^3$に対する単元小行列を選べば、
$z^4$の一意な関係が標数零からその剰余の最小多項式に移る。
各係数は$u$で高々四次なので、Vandermonde行列で平面四次式全体が決まる。
全三射影座標層で偏導関数を検査すると、幾何学的な特異点がない。
従って$t$の標数零の平面切断も滑らかである。滑らかな射影平面曲線は
Bézoutにより幾何学的に整で、超曲面自身の可約性・重複成分もこれで排除される。
超曲面の特異軌跡が二次元以上ならこの平面に会うため、特異軌跡は高々曲線である。
超曲面はCohen--Macaulayなので余次元1の正則性から正規である。アフィン錐でも
同じ判定が通る。

上で用いた基底・有限階数条件と平面切断の滑らかさは$t$の一つの開近傍で保たれる。
より詳しくは、前節の四次中心の部分束から五切断を延ばし、自由表示の剰余基底を
取る。次数16までの係数比較と有限ファイバーのtraceをこの基底で作れば、全ての
主小行列と平面判別式は正則関数となる。従ってこの開集合の各幾何点で同じ議論が成立する。

最後に$K$上で$(1,2,3,4)$の検算から第0隅の正則trace対合は非退化である。
冪零根基に属する元は全ての積のtraceを消すので、第0隅は半単純となる。
$K$では$x_1$が可逆である。各頂点の$x_1$成分は長さ4の語の和で、どの語も
第0頂点を通る。従って$e_0$は充満冪等元であり、全代数も森田同値により半単純である。
中心は$K$上四次の体なので全体は中心的単純で、中心上の次元は$1024/4=256$、
各隅は$64/4=16$となる。元の環はその中央局所化に単射であり、非零な両側イデアルも
非零に局所化するので素環である。分裂体上では十六次行列環となり、標準恒等式
$s_{32}$が成立する。閉包の標準AS点へは系\ref{cor:klein-pi-specializations}と
同じ次数別定階数の論法を使う。
\end{proof}

\begin{corollary}\label{cor:quartic-division-corners}
同じ開集合で、中心の分数体$L$上の全商環は$M_4(\Delta)$であり、
$\Delta$は$L$上次数4の中心的可除代数である。特にどの始点・終点についても、
合成可能な二つの非零経路類の積は零にならない。
\end{corollary}
\begin{proof}
非零斉次中心元を全て可逆にした中心を$L_0[u,u^{-1}]$と書ける。
ここで$L_0$は射影中心の関数体、$u$には一つの長さ4の非零中心元を取る。
従って局所化した第0隅は$\Delta_0[u,u^{-1}]$で、
$\Delta_0$はその次数0部分である。$L=L_0(u)$に延長した隅が中心的単純なので、
$\Delta_0$も$L_0$上中心的単純で次元16である。

任意の有限生成次数付き右$B$加群の$L$上の次元は64の倍数である。
実際、頂点射影$e_iB$は全て$L$上の次元が64である。
$B$は中心上有限なのでNoether環であり、標準Koszul分解から次数付き大域次元は4である。
より具体的には、四単純の左線型分解から任意の次数付き加群の
$B/B_{\ge1}$へのTorは次数4より上で消える。有限生成加群の次数付き極小分解では
従って第5項は零で、全項は有限個の頂点射影の移動の和である。階数の交代和を取ればよい。

$\Delta_0$が可除でないなら、$0,1$と異なる冪等元$v\in\Delta_0$がある。
斉次中心分母を払えば$v=b/c$で、$b\in e_0Be_0$は斉次である。
右イデアル$bB$を$L$へ延ばすと、$e_0 B_L$の真の非零直和因子$v B_L$となる。
その次元は1以上63以下であり、先の整除性に反する。よって$\Delta_0$は可除である。
可除代数上の中心変数多項式環は整域なので、$L_0(u)$への延長$\Delta$も
有限次元の整域、従って可除である。同じ議論が各隅に適用でき、
四頂点冪等元は同値な原始冪等元をなす。各$e_i B e_j$を一行一列の
$\Delta$係数に埋め込めるので、最後の積の主張も従う。
\end{proof}
